\section{Αξιολόγηση της Μεθοδολογίας}

\subsection{Datasets}
\label{Datasets}
Για το πρώτο πείραμα, που περιλαμβάνει μία μελέτη αφαίρεσης (ablation study) και μία δοκιμή των αποτελεσμάτων του επιλεγμένου μοντέλου, χρησιμοποιήθηκε το GRID Audio-Visual Speech Corpus \cite{cooke_2006_3625687}. Πρόκειται για ένα από τα πιο ευρέως χρησιμοποιούμενα dataset στον τομέα της αντίληψης ομιλίας. Αποτελείται από εγγραφές ήχου και βίντεο (της περιοχής του προσώπου) για 1000 διαφορετικές φράσεις από 34 ομιλητές (18 αρσενικούς και 16 θηλυκούς), συνολικά 34000 προτάσεις. 

Το συγκεκριμένο dataset αποτελεί παράδειγμα ενός περιορισμένου (constrained) περιβάλλοντος μάθησης, διότι οι φράσεις που καταγράφονται παρουσιάζουν μία πολύ συγκεκριμένη μορφή. Απο\-τε\-λού\-νται από ένα ρήμα εντολής (bin, lay, place, set), έπειτα ένα χρώμα (blue, green, red, white), μία πρόθεση (at, by, in, with), ένα γράμμα (οποιοδήποτε από τα 26 του αγγλικού αλφαβήτου πέραν του w), ένα ψηφίο (από το 0 έως το 9) και τέλος ένα επίρρημα (again, now, please, soon). Έτσι, το λεξιλόγιο είναι πολύ περιορισμένο σε αντίθεση με ένα in the wild dataset, στο οποίο τόσο η σύνταξη των φράσεων όσο και το ίδιο το λεξιλόγιο παρουσιάζουν σαφώς υψηλότερη ποικιλομορφία.

Για τους σκοπούς της συγκεκριμένης μελέτης, η εκπαίδευση περιορίστηκε αποκλειστικά στα 1000 βίντεο του 1ου ηθοποιού (από το s1.zip). Κάθε .mpg βίντεο έχει μήκος 3 δευτερολέπτων, ρυθμό α\-να\-νέ\-ω\-σης frames 25fps, ανάλυση 360x288 pixels και στερεοφωνικό (δικαναλικό) ήχο με ρυθμό δειγματοληψίας 44.1kHz. Η βιντεοληψία γίνεται με τον ομιλητή να κοιτάει ευθεία προς την κάμερα με ουδέτερο παρασκήνιο. Τα βίντεο έχουν τίτλους που αποδίδουν τη φράση που λέγεται σε αυτά (π.χ. η φράση ``bin green in t 5 soon'' λέγεται στο βίντεο με τίτλο bgit5s.mpg), γεγονός που βοηθάει στην εξαγωγή των πραγματικών φωνημάτων (ground truth). Πέραν των βίντεο, το dataset περιλαμβάνει αρχεία ήχου και ένα αρχείο με ακριβείς χρονισμούς λέξεων για κάθε ομιλητή, τα οποία όμως δε χρησιμοποιούνται.

\subsection{Κριτήρια Αξιολόγησης}
\label{Criteria}

Για την επακριβή αξιολόγηση της μεθοδολογίας και για τη σύγκριση με το state of the art, αποτελεί κρίσιμο βήμα η ορθή επιλογή και ο ακριβής ορισμός των μετρικών και των κριτηρίων αξιολόγησης. Σε όλα (ή κάποια) από τα παρακάτω πειράματα, η μεθοδολογία αξιολογείται με βάση τα παρακάτω κριτήρια:

\begin{itemize}

\item Γράφημα Training Loss vs Epoch: αποδίδει την τιμή της συνάρτησης κόστους στα δεδομένα του training split για κάθε epoch της εκπαίδευσης. Για τον υπολογισμό της αντικειμενικής συνάρτησης, ανατρέξτε στην ενότητα \ref{ModelStructure}.

\item Γράφημα Validation Loss vs Epoch: το ίδιο με παραπάνω αλλά στο validation split.

\item Γράφημα Validation Token Error Rate (TER) vs Epoch: ποσοστό εσφαλμένης πρόβλεψης φω\-νή\-μα\-τος σε επίπεδο token (δηλαδή, στο collapsed sequence και όχι σε επίπεδο για κάθε frame του video) στο validation split για κάθε εποχή. Υπολογίζεται ως η απόσταση edit Levenshtein διά το μήκος αναφοράς. Για πληροφορίες περί του υπολογισμού της απόστασης Levenshtein ανατρέξτε στην ενότητα \ref{OtherBackground}.

\item Γράφημα Validation Frame Accuracy vs Epoch: ποσοστό ορθών προβλέψεων φωνήματος σε επίπεδο video frame (πριν το sequence collapse) στο validation split για κάθε εποχή. Άρα, αν το βίντεο αποτελείται από $T$ frames, έχουμε 

\[
val\_frame\_acc = \frac{\# \text{ correct predictions}}{T}
\]

\item Γράφημα Training Frame Accuracy vs Epoch: το ίδιο με παραπάνω αλλά στο training split.

\item Area Under Curve (AUC): υπολογίζει την επιφάνεια κάτω από την καμπύλη (συνήθως της ακρίβειας στο validation set vs epoch), δίνοντας μία μετρική τόσο για την τελική ακρίβεια όσο και για την ταχύτητα σύγκλισης και για την αποφυγή ταλαντώσεων. 

\item Test Token Error Rate (TER): ποσοστό εσφαλμένων προβλέψεων φωνημάτων μετά το sequence collapse στο test set. Όπως παραπάνω, αξιοποιεί την απόσταση Levenshtein της πραγματικής και της προβλεπόμενης αλληλουχίας φωνημάτων.

\item Test Frame Accuracy: ποσοστό ακριβών προβλέψεων φωνημάτων σε επίπεδο frame, δηλαδή πριν το sequence collapse. 

\end{itemize}

\subsection{Πειράματα}

Παρακάτω παρουσιάζονται τα πειράματα με τα οποία δοκιμάστηκε η επίδοση της προτεινόμενης μεθοδολογίας. 

\subsubsection{Πείραμα 1: Πρόβλεψη φωνημάτων από βουβό βίντεο σε περιορισμένο περιβάλλον}
\label{Experiment1}

\paragraph{Υποπείραμα 1: Μελέτη Αφαίρεσης}

Για να απομονωθεί η συνεισφορά κάθε σχεδιαστικής επιλογής του μοντέλου, διεξήχθη μία μελέτη αφαίρεσης (ablation study) με 6 διαφορετικές εκδοχές του μοντέλου. Όλες οι εκδοχές αποτελούνται από ένα 3D CNN backbone, αξιολογήθηκαν στο ίδιο dataset με 32 κλάσεις φωνημάτων, εκπαιδεύτηκαν με AdamW optimiser με χρήση γραμμικού learning rate warmup για τις πρώτες 5 εποχές μέχρι την τιμή $3\times10^{-4}$ και scheduling ReduceLROnPlateau με τιμή patience 6 epochs και παράγοντα μείωσης 0.5. Η εκπαίδευση έγινε για μέγιστο αριθμό 40 εποχών με μηχανισμό early stopping με παράγοντα υπομονής 12 epochs στη μονάδα Ampere της Υπερυπολογιστικής Συστοιχίας ``Αριστοτέλης''. Ο παρακάτω πίνακας παρουσιάζει τις 6 διαφορετικές εκδοχές που δοκιμάστηκαν.

\begin{table}[htbp]
\centering
\caption{Οι εκδοχές του μοντέλου που δοκιμάστηκαν}
\label{tab:variants_summary}
\vspace{2mm}
\small
\begin{tabular}{|l|l|l|l|}
\hline
\textbf{Εκδοχή} & \textbf{Sequence Encoder} & \textbf{Αντικειμενική Συνάρτηση} & \textbf{Targets του CE} \\ \hline
\textbf{V1} & 2 επίπεδα BiLSTM & Frame-CE + CTC & Viseme-aware soft targets \\ \hline
\textbf{V2} & 2 επίπεδα Transformer & Frame-CE + CTC & Viseme-aware soft targets \\ \hline
\textbf{V3} & Projection head μόνο & Frame-CE + CTC & Viseme-aware soft targets \\ \hline
\textbf{V4} & 2 επίπεδα BiLSTM & CTC μόνο & N/A (χωρίς frame-CE) \\ \hline
\textbf{V5} & 2 επίπεδα BiLSTM & Frame-CE μόνο & Viseme-aware soft targets \\ \hline
\textbf{V6} & 2 επίπεδα BiLSTM & Frame-CE + CTC & Hard (identity) targets \\ \hline
\end{tabular}
\end{table}

Στη συνέχεια παρουσιάζονται τα αποτελέσματα της μελέτης, πρώτα σε μορφή πίνακα και έπειτα με μερικά χαρακτηριστικά γραφήματα.

Να σημειωθεί ότι οι απόλυτες τιμές των τιμών της αντικειμενικής συνάρτησης στα σχήματα \ref{fig:ablation_train_loss}, \ref{fig:ablation_train_val_loss_grid} και \ref{fig:ablation_val_loss} δεν είναι συγκρίσιμες. Συγκεκριμένα, τα μοντέλα με μόνο έναν από τους 2 όρους (είτε μόνο cross entropy είτε μόνο CTC) έχουν αντικειμενικές συναρτήσεις με μικρότερες τιμές σε απόλυτη τιμή, καθώς ο πολλαπλασιαστικός συντελεστής των δύο όρων παραμένει σταθερός. Τα σχήματα στο \ref{fig:ablation_val_loss_components} ωστόσο, που παρουσιάζουν τη συνεισφορά κάθε όρου της αντικειμενικής συνάρτησης χωριστά, είναι απολύτως συγκρίσιμα. Σημαντικό, επίσης, είναι να παρατηρηθεί ότι το μοντέλο V6 παρουσιάζει τη χαμηλότερη τιμή training και validation loss ανάμεσα στα V1, V2, V3 και V6, χωρίς ωστόσο να πετυχαίνει και τη βέλτιστη επίδοση όσον αφορά το test TER. Αυτό οφείλεται στο ότι η χρήση soft targets στα μοντέλα V1, V2 και V3 μοιράζει την πιθανότητα κάθε φωνήματος σε διάφορες όμοιες (οπτικά ή ακουστικά) κλάσεις παράγοντας ένα target με υψηλότερη εντροπία και άρα οδηγώντας σε υηψλότερες τιμές cross entropy, σε σχέση με το one-hot encoding που χρησιμοποιείται για το μοντέλο V6.

\begin{table}[htbp]
\centering
\caption{Περίληψη των αποτελεσμάτων της μελέτης αφαίρεσης}
\label{tab:ablation_summary}
\vspace{2mm}
\small
\begin{tabular}{|p{5.2cm}|c|c|c|c|c|}
\hline
\textbf{Variant} & \textbf{\begin{tabular}[c]{@{}c@{}}Best Val\\ TER\end{tabular}} & \textbf{\begin{tabular}[c]{@{}c@{}}Test\\ TER\end{tabular}} & \textbf{\begin{tabular}[c]{@{}c@{}}Test Frame\\ Acc.\end{tabular}} & \textbf{\begin{tabular}[c]{@{}c@{}}AUC\\ (Val $1-\text{TER}$)\end{tabular}} & \textbf{\begin{tabular}[c]{@{}c@{}}Εποχές μέχρι\\ Val TER $<0.20$ \end{tabular}} \\ \hline
V1 — BiLSTM + Full Loss & \textbf{0.1251} & \textbf{0.1219} & \textbf{0.8474} & \textbf{0.7840} & \textbf{8} \\ \hline
V2 — Transformer + Full Loss & 0.1551 & 0.1566 & 0.8394 & 0.7176 & 16 \\ \hline
V3 — Projection head μόνο & 0.2540 & 0.2464 & 0.8122 & 0.6006 & Ποτέ \\ \hline
V4 — BiLSTM + CTC Only & 0.1977 & 0.1869 & 0.4743 & 0.6005 & 33 \\ \hline
V5 — BiLSTM + Frame-CE Μόνο & 0.1555 & 0.1419 & 0.8401 & 0.7402 & 16 \\ \hline
V6 — BiLSTM + Hard Targets & 0.1266 & 0.1289 & 0.8456 & 0.7696 & 9 \\ \hline
\end{tabular}
\end{table}

\begin{figure}[H]
	\centering
	\includegraphics[width=0.7\textwidth]{ablation_train_loss.png}
	\caption{Γράφημα τιμών αντικειμενικής συνάρτησης στο training set για κάθε epoch}
	\label{fig:ablation_train_loss}
\end{figure}

\begin{figure}[H]
	\centering
	\includegraphics[width=0.7\textwidth]{ablation_train_frame_acc.png}
	\caption{Γράφημα τιμών frame accuracy στο training set για κάθε epoch}
\end{figure}

\begin{figure}[H]
	\centering
	\includegraphics[width=0.7\textwidth]{ablation_train_val_loss_grid.png}
	\caption{Σύγκριση τιμών loss για training και validation κάθε μοντέλου}
	\label{fig:ablation_train_val_loss_grid}
\end{figure}

\begin{figure}[H]
	\centering
	\includegraphics[width=0.7\textwidth]{ablation_val_loss.png}
	\caption{Γράφημα τιμών αντικειμενικής συνάρτησης στο validation set για κάθε epoch}
	\label{fig:ablation_val_loss}
\end{figure}

\begin{figure}[H]
	\centering
	\includegraphics[width=0.7\textwidth]{ablation_val_loss_components.png}
	\caption{Ανάλυση συνεισφοράς των όρων της αντικειμενικής συνάρτησης (CE και CTC)}
	\label{fig:ablation_val_loss_components}
\end{figure}

\begin{figure}[H]
	\centering
	\includegraphics[width=0.7\textwidth]{ablation_val_frame_acc.png}
	\caption{Γράφημα τιμών frame accuracy στο validation set για κάθε epoch}
\end{figure}

\begin{figure}[H]
	\centering
	\includegraphics[width=0.7\textwidth]{ablation_val_ter.png}
	\caption{Γράφημα τιμών TER στο validation set για κάθε epoch}
\end{figure}

\begin{figure}[H]
	\centering
	\includegraphics[width=0.7\textwidth]{ablation_lr_schedule.png}
	\caption{Γράφημα τιμών του ρυθμού μάθησης (learning rate) για κάθε epoch}
\end{figure}

Από τα παραπάνω εξάγεται το συμπέρασμα ότι το μοντέλο της μεθοδολογίας παρουσιάζει την ισχυρότερη επίδοση τόσο όσον αφορά την τελική απόδοση όσο και την ομαλότητα και αποδοτικότητα της εκπαίδευσης, γεγονότα που αποδεικνύονται από την υπεροχή του V1 στις μετρικές test TER, test frame accuracy και AUC. Στις επόμενες παραγράφους δίνεται μία αναλυτική σύγκριση κάθε εκδοχής με την V1, όπως και τα απαραίτητα σχόλια.

Όπως αποδεικνύει η σύγκριση μεταξύ V1 και V3, η σχεδιαστική επιλογή με το υψηλότερο αντίκρισμα στα αποτελέσματα είναι η παρουσία του χρονικού sequence encoder. Εν απουσία του BiLSTM και αξιοποιώντας μόνο ένα γραμμικό projection head, το μοντέλο V3 πετυχαίνει τελικά test TER 0.2464, υπερδιπλασιάζοντας αυτό του μοντέλου V1. Οι υπόλοιπες μετρικές (val TER, val frame accuracy, AUC) αποδεικνύουν επίσης την αναγκαιότητα χρήσης sequence encoder. Να σημειωθεί, ωστόσο, ότι μετρικές όπως το test frame accuracy, με τιμή 0.812, αποδεικνύουν ότι το μοντέλο πράγματι είναι ικανό να μάθει. Περιορίζεται, ωστόσο, από την έλλειψη της πληροφορίας παρακειμένου (contextual information), το οποίο είναι βασικό στοιχείο για την κατανόηση των φωνημάτων. Το context που προσφέρει η χρήση 3D CNN backbone δεν είναι τόσο πλούσιο όσο αυτό του συνδυασμού 3D CNN και BiLSTM.

Αντικαθιστώντας τα επίπεδα BiLSTM με Transformer ίδιου βάθους (V2), παρατηρούμε μία σημαντική επιδείνωση των αποτελεσμάτων, της τάξης περίπου του 25\%. Η πορεία της εκπαίδευσης δείχνει μία πιο αργή βελτίωση του val TER στο V2 σε σχέση με το V1. Ένα ακόμα στοιχείο που αξίζει να σημειωθεί είναι ότι ο μηχανισμός scheduling του ρυθμού μάθησης ReduceLROnPlateau δε πυροδοτήθηκε κατά την εκπαίδευση του V2, ενώ προκλήθηκε 2 φορές κατά την εκπαίδευση του V1, αποδεικνύοντας ότι η σύγκλιση του V2 δεν έχει ακόμα ολοκληρωθεί κατά την 40η εποχή και ότι θα χρειαζόταν υψηλότερο υπολογιστικό κόστος εκπαίδευσης για να παρουσιάσει το V2 παρόμοια αποτελέσματα με αυτά του V1. Μία από τις κύριες αιτίες για τα αποτελέσματα του V2 είναι το μέγεθος του dataset που χρησιμοποιήθηκε, το οποίο περιλαμβάνει μόλις 700 φράσεις. Οι transformer encoders απαιτούν περισσότερα και πιο ποικιλόμορφα δεδομένα για να μάθουν αποδοτικά μοτίβα attention. Αντιθέτως, το inductive bias των BiLSTM για χρονική εξάρτηση, ειδικά για γειτονικό context, λόγω της recurrent αρχιτεκτονικής τους, βρίσκει εύκολα εφαρμογή στην κίνηση των χειλιών, κατά την οποία γειτονικά frames αλληλεξαρτώνται σε υψηλό βαθμό.

Συγκρίνοντας τα μοντέλα V1 και V4, αναδεικνύεται η σημασία του όρου cross entropy της αντικειμενικής συνάρτησης. Το μοντέλο V4 πετυχαίνει στη μετρική test TER τιμή 0.1869, πάνω από 50\% επιδεινωμένη επίδοση σε σχέση με το V1. Η διαφορά των δύο μοντέλων γίνεται ιδιαίτερα προφανής με τη μετρική test frame accuracy, η οποία για το V4 είναι μόλις 0.4743, τιμή που ανταποκρίνεται με την τυχαία ανάθεση labels σε πρόβλημα ταξινόμησης με μία σαφώς επικρατέστερη κλάση, αυτήν του blank token (να μην εκφέρεται στο δεδομένο frame κάποιο φώνημα). Τα αποτελέσματα αυτά είναι αναμενόμενα: το σήμα του gradient, παραγόμενο αποκλειστικά μέσω του όρου CTC, οδηγεί το μοντέλο να μάθει να προβλέπει σωστές αλληλουχίες φωνημάτων αλλά όχι ποιες κινήσεις των χειλιών (και άρα ποια frames του βίντεο) παράγουν το οποιοδήποτε φώνημα. Έτσι, η διαφορά στο test TER δεν είναι τόσο μεγάλη όσο στο test frame accuracy.

Από την άλλη, αφαιρώντας τον όρο CTC και διατηρώντας μόνο τον όρο cross entropy της συνάρτησης κόστους (με το μοντέλο V5), επιτυγχάνεται τελικά test TER με τιμή 0.1419, ελαφρώς χειρότερα από το μοντέλο V1, τιμή που αποτελεί τη 3η καλύτερη στη συγκεκριμένη μετρική. Το σημαντικότερο πρόβλημα που αντιμετωπίζει η εκδοχή V5 είναι η ομαλότητα της εκπαίδευσης. Συγκεκριμένα, η καμπύλη validation TER vs epoch παρουσιάζει εμφανείς ταλαντώσεις και η σύγκλιση είναι πιο αργή σε σχέση με τα μοντέλα V1, V2 και V6. Ταυτόχρονα, το γράφημα val loss vs epoch παρουσιάζει στοιχεία overfitting, καθώς η τιμή αυξάνει μετά την 20η εποχή. Έτσι, εξάγεται το συμπέρασμα ότι ο όρος CTC συμβάλλει στην κανονικοποίηση του μοντέλου, προσθέτοντας τη συνθήκη ότι το τελικό σύνολο προβλέψεων φωνημάτων ανά frame θα πρέπει να καταρρέει σε μία λογική αλληλουχία φωνημάτων. Με αυτόν τον τρόπο βελτιώνεται η γενίκευση αλλά και η απόδοση του μοντέλου.

Τέλος, η διαφορά των μοντέλων V1 και V6 έγκειται στα targets του cross entropy της αντικειμενικής συνάρτησης: ενώ το V1 χρησιμοποιεί soft targets με βάση το πόσο όμοια είναι τα φωνήματα, ακουστικά ή οπτικά, το V6 χρησιμοποιεί τα συνήθη hard targets του cross entropy. Το μοντέλο V6 πετυχαίνει τη 2η καλύτερη τιμή test TER (0.1289), ενώ παρουσιάζει εικόνα εκπαίδευσης πολύ παρόμοια με αυτή του μοντέλου V1. Η ελάχιστη διαφορά των δύο εκδοχών φαίνεται να σημαίνει ότι η χρήση των soft targets παρέχει μικρό αλλά ίσως σταδιακά αυξανόμενο πλεονέκτημα με βάση το μέγεθος και την ποικιλομορφία του dataset. Όσο αυξάνεται η συχνότητα παρόμοιων φωνημάτων, τόσο ο μηχανισμός αυτός θα παρείχε περισσότερη ομαλότητα κατά την εκπαίδευση, προστατεύοντας το μοντέλο από θόρυβο labelling σε αμφίβολες περιπτώσεις. Στο μικρό και περιορισμένο dataset που χρησιμοποιείται στο πείραμα αυτό, είναι πιο εύκολο να "αποστηθίσει" το μοντέλο αλληλουχίες φωνημάτων (άρα λέξεων, π.χ. να μάθει ότι το φώνημα /b/ βρίσκεται πριν το φώνημα /i/ για να δημιουργηθεί η λέξη "bin", και όχι τα οπτικά πανομοιότυπα φωνήματα /p/ ή /m/), σε αντίθεση με ένα ευρύτερο in the wild dataset (στο οποίο κάλλιστα θα μπορούσε να περιέχεται και η λέξη "pin" και άρα να πρέπει να προβλεφθεί το /p/ και όχι το /b/). 

\paragraph{Υποπείραμα 2: Δοκιμή του επιλεγμένου μοντέλου σε περιορισμένο περιβάλλον}

Έχοντας καταλήξει στην τελική και βέλτιστη μορφή του μοντέλου και έχοντας επεξηγήσει τη συνεισφορά κάθε σχεδιαστικής επιλογής στο τελικό αποτέλεσμα, ο πειραματισμός συνεχίζεται με μία δοκιμή σε μεγαλύτερο βάθος στο ίδιο dataset (GRID Audio-Visual Speech Corpus). Συγκεκριμένα, πέραν της πρόβλεψης των φωνημάτων, η ανάλυση προχωρά και σε πρόβλεψη λέξεων με γνώση του περιορισμένου περιβάλλοντος της μελέτης, προσομοιώνοντας έτσι ένα σύστημα με πιο πραγματική εφαρμογή.

Η εκπαίδευση έγινε πάλι με AdamW optimiser με χρήση γραμμικού learning rate warmup για τις πρώτες 5 εποχές μέχρι την τιμή $3\times10^{-4}$ και scheduling ReduceLROnPlateau με τιμή patience 6 epochs και παράγοντα μείωσης 0.5, στη μονάδα Ampere της Υπερυπολογιστικής Συστοιχίας ``Αριστοτέλης''. Ο μέγιστος αριθμός εποχών αυξήθηκε σε 60, διατηρώντας τον μηχανισμό early stopping με παράγοντα υπομονής 12 epochs. Παρακάτω παρουσιάζονται τα αποτελέσματα του συγκεκριμένου πειράματος:

\begin{figure}[H]
	\centering
	\includegraphics[width=0.7\textwidth]{constrained_train_losses.png}
	\caption{Γράφημα τιμών του train loss και του κάθε συντελεστή αυτού για κάθε epoch}
	\label{fig:constrained_train_losses}
\end{figure}

\begin{figure}[H]
	\centering
	\includegraphics[width=0.7\textwidth]{constrained_val_loss.png}
	\caption{Γράφημα τιμών του validation loss για κάθε epoch}
\end{figure}

\begin{figure}[H]
	\centering
	\includegraphics[width=0.7\textwidth]{constrained_val_ter.png}
	\caption{Γράφημα τιμών του validation TER για κάθε epoch}
\end{figure}

\begin{figure}[H]
	\centering
	\includegraphics[width=0.7\textwidth]{constrained_lr_schedule.png}
	\caption{Γράφημα τιμών του ρυθμού μάθησης (learning rate) για κάθε epoch}
\end{figure}

Αρχικά, να σημειωθεί ότι στο σχήμα \ref{fig:constrained_train_losses}, ο συντελεστής CTC της συνάρτησης κόστους παρουσιάζεται πριν τον πολλαπλασιασμό του με τον κατάλληλο συντελεστή, άρα για κάθε epoch (όπως είναι προφανές από το σχήμα), δεν ισχύει $total\_loss = frame\_ce\_loss + ctc\_loss$, αλλά $total\_loss = frame\_ce\_loss + 0.3 \times ctc\_loss$. 

Παρατηρείται ομαλή εκπαίδευση που τελικά οδηγεί σε βέλτιστη τιμή validation loss 1.027 όπως και βέλτιστη τιμή validation TER 0.1273, με τον μηχανισμό early stopping να πυροδοτείται κατά την εποχή 50. Όπως ήταν και αναμενόμενο, τα αποτελέσματα είναι πολύ παρόμοια με αυτά του μοντέλου V1 της μελέτης αφαίρεσης. Οι μικροδιαφορές οφείλονται σε ελάχιστες διαφοροποιήσεις κατά την εκτέλεση των εκπαιδεύσεων και όχι σε κάποια αρχιτεκτονική παράλειψη. 

Στο test set, το μοντέλο πετυχαίνει τιμή $frame\_ce\_loss = 0.8626$, $ctc\_loss = 0.461$, $total\_loss = 1.0009$ και $ter =  0.1241$. Παρακάτω παρουσιάζεται ο πίνακας σύγχυσης (confusion matrix) των φωνημάτων που προβλέφθηκαν από το μοντέλο σε σχέση με τα πραγματικά (ground truth). 

\begin{figure}[H]
	\centering
	\includegraphics[width=0.7\textwidth]{constrained_confusion_matrix.png}
	\caption{Πίνακας σύγχυσης (confusion matrix) στο test set}
\end{figure}

Ο πίνακας σύγχυσης παρουσιάζει ισχυρή προβλεπτική ικανότητα, με την κύρια διαγώνιο (το σύνολο των σωστών προβλέψεων) να κυριαρχεί σε σχέση με τον υπόλοιπο πίνακα. Οι κύριες πηγές προβληματισμού είναι η τελευταία σειρά ``<inserted>'', η οποία αποτελείται από περιπτώσεις που ένα φώνημα προβλέφθηκε λανθασμένα στη θέση του blank token και η τελευταία στήλη ``<missed>'' που περιγράφει το αντίστροφο φαινόμενο, όταν δηλαδή δεν προβλέπεται κάποιο φώνημα που όντως υπάρχει αλλά στη θέση του προβλέπεται το blank token. Συγχύσεις όμοιων οπτικά φωνημάτων (π.χ. /d/ και /t/) υπάρχουν αλλά σε ελάχιστο βαθμό, γεγονός που εξηγείται από την περιορισμένη φύση του dataset: το μοντέλο μαθαίνει ότι το dataset αποτελείται από συγκεκριμένες λέξεις και προσαρμόζεται ανάλογα.

Σε συνέχεια των παραπάνω και εκμεταλλευόμενοι τον περιορισμό των λέξεων που χρησιμοποιούνται στο dataset, αναπτύχθηκε ένα script πρόβλεψης φράσεων από την αλληλουχία φωνημάτων. Το script αυτό έχει πλήρη γνώση τόσο του λεξιλογίου του dataset όσο και της γραμματικής που παρουσιάζουν οι φράσεις (θεωρείται γνωστό, δηλαδή, ότι οι φράσεις είναι της μορφής εντολή, χρώμα, πρόθεση, γράμμα, ψηφίο, επίρρημα, όπως και ποιες είναι οι πιθανές λέξεις για κάθε μία από αυτές της κατηγορίες). Με βάση την προφορά του ηθοποιού, δημιουργείται ένα λεξικό που αντιστοιχίζει αλληλουχίες φωνημάτων σε λέξεις. Ο αλγόριθμος που αναπτύχθηκε περιγράφεται με περρισότερη λεπτομέρεια στην ενότητα (ρεφ).

Η χρήση του παραπάνω script στις προβλέψεις του βέλτιστου μοντέλου πετυχαίνει ορθή πρόβλεψη σε επίπεδο λέξης για 87.1\% των λέξεων του test set και ακριβή πρόβλεψη σε επίπεδο πρότασης σε ποσοστό 38\%. Το κύριο πρόβλημα είναι η πρόβλεψη στην ομάδα των γραμμάτων: μόλις 46\% των γραμμάτων προβλέπονται σωστά. Αυτό οφείλεται τόσο στο μέγεθος της ομάδας αυτής (25 υποψήφιες λέξεις έναντι 4 σε άλλες κατηγορίες), όπως και στο μέγεθος της αλληλουχίας των φωνημάτων που αποτελούν ένα γράμμα (συνήθως είναι λιγότερα και άρα η πρόβλεψη είναι πιο επιρρεπής στον θόρυβο). Η υπόθεση αυτή ενισχύεται από την κατηγορία των προθέσεων, η οποία εμφανίζει το δεύτερο χαμηλότερο ποσοστό ορθής πρόβλεψης με 85.3\%. Αντιθέτως, στις ομάδες με μεγαλύτερες λέξεις ή με λέξεις που απαρτίζονται από πιο χαρακτηριστικά φωνήματα, η πρόβλεψη ήταν ιδιαίτερα ακριβής: 98\% στις εντολές, 99.3\% στα χρώματα, 95.3\% στα ψηφία και 98.7\% στα επιρρήματα. Τα ποσοστά αυτά είναι εξαιρετικά και αποδεικνύουν ότι οι προβλέψεις τόσο σε επίπεδο φωνήματος από το μοντέλο όσο και σε επίπεδο λέξης από τον αλγόριθμο είναι ιδιαίτερα ικανές στο περιορισμένο περιβάλλον.

Για να κριθεί ο χρονισμός (timing) της πρόβλεψης και όχι μόνο η ορθότητα του περιεχομένου, χρησιμοποιήθηκε το ffmpeg ώστε να υποτιτλιστούν μερικά από τα video του test set με τις λέξεις που προβλέπονται στον χρόνο που προβλέπονται τα φωνήματα τα οποία παράγουν τις λέξεις. Η διαδικασία αυτή δεν μπορεί να παρουσιαστεί στο πλαίσιο της συγκεκριμένης αναφοράς, παρουσιάζει ωστόσο ιδιαίτερα ακριβή αποτελέσματα, ειδικά, φυσικά, όταν η πρόβλεψη ολόκληρης της φράσης είναι σωστή. 
\subsubsection{Πείραμα 2: Πρόβλεψη φωνημάτων από βουβό βίντεο σε μη περιορισμένο (in the wild) περιβάλλον}

\subsubsection{Πείραμα 3: Πρόβλεψη φωνημάτων με συνδυασμό οπτικών και ακουστικών στοιχείων}
\label{Experiment 3}

Στο συγκεκριμένο πείραμα δοκιμάζεται η μεθοδολογία του audiovisual μοντέλου που παρουσιάζεται \ref{AudiovisualModel}, χρησιμοποιώντας δηλαδή ως είσοδο στο μοντέλο όχι μόνο το βουβό βίντεο, αλλά και το τμήμα του ήχου που αντιστοιχεί σε αυτό (ακριβέστερα, το παραγόμενο Mel-Spectrogram). Για τον σκοπό αυτό, χρησιμοποιούνται τόσο κατά την εκπαίδευση με το training set όσο και κατά τη δοκιμή με τα validation και testing sets δεδομένα στα οποία έχει προστεθεί θόρυβος με τους τρόπους που παρουσιάστηκαν παραπάνω, στην ενότητα \ref{AudiovisualPreprocessing}.

\paragraph{Υποπείραμα 1: Αύξηση του SNR του ήχου}
\label{SNRExperiment}

Κατά το 1ο υποπείραμα, μελετάται η επίδραση της μείωσης του ηχητικού θορύβου στην ικανότητα του audiovisual μοντέλου να προβλέπει τα σωστά φωνήματα, ενώ συγκρίνεται και με το μοντέλο το οποίο χρησιμοποιεί αποκλειστικά το ηχητικό σήμα για την πρόβλεψη. Το αναμενόμενο αποτέλεσμα είναι ότι, με την αύξηση του SNR, το ηχητικό σήμα γίνεται όλο και πιο αξιόπιστο, με αποτέλεσμα το μοντέλο να ``εμπιστεύεται'' τον ήχο περισσότερο και να προβλέπει με βάση αυτό το σήμα με μεγαλύτερη ακρίβεια. Παράλληλα, αναμένουμε ότι το μοντέλο που χρησιμοποιεί και τα δύο modalities θα έχει βελτιωμένη επίδοση σε σχέση με το αντίστοιχο που χρησιμοποιεί αποκλειστικά τον ήχο. Για τη συγκεκριμένη δοκιμή, τα frames των βίντεο που επιλέγονται για να τους προστεθεί θόρυβος (με τη μορφή του μηδενισμού, του παγώματος ή του θολώματος) είναι κατά μέσο όρο 30 (από τα συνολικά 75 του κάθε βίντεο του GRID), με jitter $\pm$ 8 frames. Ο θόρυβος που προστίθεται στον ήχο χαρακτηρίζεται από SNR στις τιμές -5, 0, 5, 10 και 15, ενώ χρησιμοποιείται και dropout του ηχητικού σήματος με πιθανότητα 30\% και μήκος 15 frames (για κάθε διαφορετική φράση, μέχρι 1 φορά μπορεί να γίνει dropout).

Κατά την εκπαίδευση όλων των μοντέλων, χρησιμοποιήθηκε Adam optimizer με ρυθμό μάθησης 3e-4, weight decay με τιμή 1e-4 και gradient clipping με τιμή 1. Χρησιμοποιήθηκε scheduling του ρυθμού μάθησης με γραμμικό warmup για τις πρώτες 5 εποχές, μείωση του ρυθμού μάθησης μετά από 6 εποχές χωρίς βελτίωση κατά παράγοντα 0.5, με ελάχιστη τιμή το 1e-6.  Χρησιμοποιείται early stopping με κριτήριο τη μετρική val TER (validation Token Error Rate) και τιμή υπομονής 12 εποχές, batch size 4, με μέγιστο αριθμό εποχών να είναι 100. Τέλος, κατά την εκπαίδευση, χρησιμοποιείται παράγοντας modality dropout 0.15 (με πιθανότητα 0.15 αφαιρείται πλήρως ένα από τα δύο modalities) για να αυξηθεί η ευρωστία του μοντέλου. Μέσω των παρακάτω διαγραμμάτων, αποτυπώνονται τα αποτελέσματα του συγκεκριμένου υποπειράματος.

Όπως φαίνεται στο σχήμα \ref{best_val_ter_vs_snr_db}, η μετρική val TER (Token Error Rate στο validation set) βελτιώνεται τόσο για το audiovisual μοντέλο, όσο και για το audio unimodal μοντέλο με την αύξηση του SNR του ηχητικού σήματος (με εξαίρεση μία μικρή αύξηση του TER για 15dB έναντι των 10dB στην περίπτωση του audiovisual μοντέλου). Το φαινόμενο αυτό είναι αναμενόμενο και φαίνεται να επιβεβαιώνει την ικανότητα του μοντέλου να αξιοποιήσει περισσότερο την πληροφορία του ήχου, όσο αυτή γίνεται πιο αξιόπιστη. Ταυτόχρονα, τα μοντέλα που αξιοποιούν και τα δύο modalities (βίντεο και ήχο) υπερισχύουν σε κάθε περίπτωση από τα αντίστοιχα που αξιοποιούν μονάχα το ηχητικό σήμα. Αυτό (σε συνδυασμό με το σχήμα \ref{best_val_ter_vs_frames} παρακάτω) αποδεικνύει πως σε περιβάλλοντα όπου υπάρχει θόρυβος και στο βίντεο και στον ήχο, είναι προτιμότερο να λαμβάνονται υπόψιν και οι 2 μορφές δεδομένων ώστε να αναγνωρίζονται με υψηλότερο ποσοστό επιτυχίας τα φωνήματα και τελικά η ανθρώπινη ομιλία.

\begin{figure}[H]
	\centering
	\includegraphics[width=0.7\textwidth]{best_val_ter_vs_snr_db.png}
	\caption{Γραφήματα τιμών του βέλτιστου val TER για διάφορες τιμές audio SNR}
	\label{best_val_ter_vs_snr_db}
\end{figure}

Στο σχήμα \ref{epoch_train_loss_across_snr_db}, παρουσιάζονται οι καμπύλες εκπαίδευσης (train loss) των μοντέλων για διάφορες τιμές SNR του ηχητικού σήματος. Να σημειωθεί ότι οι τιμές της αντικειμενικής συνάρτησης των audiovisual και των αντίστοιχων audio μοντέλων δεν είναι απολύτως συγκρίσιμες, καθώς η αντικειμενική συνάρτηση του audiovisual μοντέλου περιλαμβάνει έναν περαιτέρω auxiliary CTC όρο, συγκεκριμένα αυτόν του βίντεο. Γενικά, φαίνεται σε όλα τα μοντέλα να επιτυγχάνεται σταθερή σύγκλιση, με την τιμή του SNR να μην επηρεάζει την εικόνα των καμπυλών. Με την αύξηση του SNR, φαίνεται να βελτιώνεται ελάχιστα η τελική τιμή της αντικειμενικής συνάρτησης στο training set. 

\begin{figure}[H]
	\centering
	\includegraphics[width=0.9\textwidth]{epoch_train_total_by_snr_db.png}
	\caption{Γραφήματα τιμών train loss per epoch για διάφορες τιμές audio SNR}
	\label{epoch_train_loss_across_snr_db}
\end{figure}

Τα σχήματα \ref{epoch_audio_gate_across_snr_db} και \ref{epoch_video_gate_across_snr_db} δείχνουν τις μέσες τιμές της audio gate και της video gate, αντίστοιχα, για κάθε εποχή της εκπαίδευσης. Προφανώς, στην περίπτωση των audio μοντέλων, οι τιμές αυτές είναι 1 και 0 αντίστοιχα για κάθε εποχή, καθώς δε χρησιμοποιείται εξ' ορισμού το βίντεο. Παρουσιάζει ιδιαίτερο ενδιαφέρον, ωστόσο, η μελέτη των αποτελεσμάτων για το audiovisual μοντέλο. Σε γενικές γραμμές, αναμένουμε το μοντέλο να μαθαίνει να εμπιστεύεται τον ήχο όλο και περισσότερο, όσο το SNR αυξάνεται, επομένως το audio gate να αυξάνει με το SNR και αντίθετα το video gate να μειώνεται. Πράγματι, η τελική εικόνα φαίνεται να επιβεβαιώνει αυτό το φαινόμενο. Παρατηρούμε ότι, για τιμή SNR -5dB, όπου δηλαδή ο θόρυβος υπερισχύει του πραγματικού σήματος, η τιμή της audio gate είναι σταθερά μικρότερη από αυτήν της video gate, δηλαδή ότι το μοντέλο, ορθά, βασίζει τις αποφάσεις του κυρίως στο βίντεο και όχι τόσο στον ήχο. Για SNR 0dB, η τιμή της audio gate φαίνεται να ταλαντεύεται γύρω από το 0.5, υποδεικνύοντας ότι το μοντέλο εμπιστεύεται το ηχητικό σήμα και το βίντεο με όμοιο βάρος. Ήδη, δηλαδή, από το σημείο που ο πραγματικός ήχος και ο θόρυβος έχουν ίδια ενέργεια, το ηχητικό σήμα γίνεται εξίσου αξιόπιστο με το βίντεο από το οποίο έχουν αλλαχθεί λιγότερα από τα μισά frames του. Αυτό αποδεικνύει ότι, σε γενικές γραμμές, με τη βοήθεια του ηχητικού σήματος εξάγονται πιο εύκολα συμπεράσματα για τα φωνήματα και τελικά για την ομιλία, πράγμα αναμενόμενο από όσα έχουν ήδη συζητηθεί. Για τις υπόλοιπες τιμές του SNR, το μοντέλο ορθά εμπιστεύεται όλο και περισσότερο το ηχητικό σήμα εις βάρος του βίντεο για να εξάγει την τελική ταξινόμηση, με τιμή audio gate π.χ. για τα 15dB να ξεπερνάει στο τέλος της εκπαίδευσης το 0.6. Τα αποτελέσματα αυτά φαίνεται να επιβεβαιώνουν την ικανότητα του μοντέλου να προσαρμοστεί στις συνθήκες στις οποίες εφαρμόζεται και να αξιοποιεί την εκάστοτε modality που μπορεί πράγματι να οδηγήσει στη σωστή ταξινόμηση, παρουσιάζοντας μία γενική προτίμηση για το ηχητικό σήμα, κάτι που είναι λογικό, μια και ο ήχος μπορεί απόλυτα να χαρακτηρίσει ένα φώνημα, ενώ το βίντεο όχι.

\begin{figure}[H]
	\centering
	\includegraphics[width=0.9\textwidth]{epoch_val_gate_audio_by_snr_db.png}
	\caption{Γραφήματα τιμών της audio gate στο validation set per epoch για διάφορες τιμές audio SNR}
	\label{epoch_audio_gate_across_snr_db}
\end{figure}

\begin{figure}[H]
	\centering
	\includegraphics[width=0.9\textwidth]{epoch_val_gate_video_by_snr_db.png}
	\caption{Γραφήματα τιμών της video gate στο validation set per epoch για διάφορες τιμές audio SNR}
	\label{epoch_video_gate_across_snr_db}
\end{figure}

Τα σχήματα \ref{epoch_val_ter_across_snr_db} και \ref{epoch_val_loss_across_snr_db} δείχνουν την πορεία της μετρικής val TER και της loss function στο validation set για τις εποχές της εκπαίδευσης, για τις διάφορες τιμές του SNR. Να σημειωθεί ότι και σε αυτήν την περίπτωση οι τιμές της αντικειμενικής συνάρτησης των audiovisual μοντέλων και των audio μοντέλων δεν είναι απολύτως συγκρίσιμες, για τον ίδιο λόγο που αναφέρθηκε παραπάνω. Όσον αφορά τη μετρική val TER, παρατηρούμε σύγκλιση για όλα τα μοντέλα, με χαμηλότερη τελική τιμή για το audiovisual μοντέλο εις βάρος του αντίστοιχου audio μοντέλου σε κάθε περίπτωση, όπως ήταν και αναμενόμενο. Τόσο η ταχύτητα και η ομαλότητα της σύγκλισης όσο και η τελική τιμή του val TER φαίνεται να βελτιώνονται με την αύξηση του SNR του ηχητικού σήματος. Τα διαγράμματα των τιμών της αντικειμενικής συνάρτησης στο validation set παρουσιάζουν όμοια εικόνα με τα αντίστοιχα στο training set του σχήματος \ref{epoch_train_loss_across_snr_db}, με καθολική ομαλή σύγκλιση. Η τελική τιμή της validation loss φαίνεται να βελτιώνεται όσο αυξάνεται το SNR, τόσο στο audiovisual όσο και στο audio μοντέλο.

\begin{figure}[H]
	\centering
	\includegraphics[width=0.9\textwidth]{epoch_val_ter_by_snr_db.png}
	\caption{Γραφήματα τιμών val TER per epoch για διάφορες τιμές audio SNR}
	\label{epoch_val_ter_across_snr_db}
\end{figure}

\begin{figure}[H]
	\centering
	\includegraphics[width=0.9\textwidth]{epoch_val_total_by_snr_db.png}
	\caption{Γραφήματα τιμών val loss per epoch για διάφορες τιμές audio SNR}
	\label{epoch_val_loss_across_snr_db}
\end{figure}

Το σχήμα \ref{epochs_vs_snr_db} δείχνει τις εποχές πριν πυροδοτηθεί ο μηχανισμός early stopping για το audiovisual και το audio μοντέλο ως συνάρτηση του SNR του ήχου που χρησιμοποιείται κατά την εκπαίδευση. Η γενική τάση φαίνεται να είναι το audiovisual μοντέλο να απαιτεί περισσότερες εποχές εκπαίδευσης μέχρι τη σύγκλιση, κάτι που πιθανότατα εξηγείται από την ύπαρξη και των δύο modalities. Σε αντίθεση με το audio μοντέλο, το audiovisual μοντέλο χρησιμοποιεί και οπτικά και ακουστικά στοιχεία για την εξαγωγή των τελικών του αποφάσεων, με αποτέλεσμα να δαπανεί περισσότερες εποχές ώστε να αντιληφθεί σε ποια σημεία πρέπει να εμπιστεύεται το βίντεο και σε ποια τον ήχο. Παράλληλα, μαθαίνει να συνδυάζει τα συμπεράσματα του ενός modality με αυτά του άλλου, κάτι που το audio μοντέλο δε χρειάζεται να κάνει. Η τάση αυτή δε συνεχίζεται όταν το SNR του ήχου είναι 15dB. Σε εκείνο το σημείο, το μοντέλο αξιοποιεί σημαντικά περισσότερο τον ήχο (όπως φαίνεται στο σχήμα \ref{epoch_audio_gate_across_snr_db}), βασιζόμενο στο βίντεο μόνο όταν αυτό είναι απαραίτητο (π.χ. στα σημεία με πλήρες audio dropout). Πρόκειται, δηλαδή, για το σημείο που το audiovisual μοντέλο μπορεί βέλτιστα να συνδυάσει τις πληροφορίες των δύο modalities και να συγκλίνει ταχύτερα σε σχέση με το αντίστοιχο audio μοντέλο.

\begin{figure}[H]
	\centering
	\includegraphics[width=0.7\textwidth]{epochs_ran_vs_snr_db.png}
	\caption{Γράφημα αριθμού εποχών εκπαίδευσης για διάφορες τιμές audio SNR}
	\label{epochs_vs_snr_db}
\end{figure}

Τα διαγράμματα \ref{test_audio_gate_vs_db} και \ref{test_video_gate_vs_db} δείχνουν τις μέσες τιμές της audio gate και της video gate αντίστοιχα στο test set. Φυσικά, στην περίπτωση του audio μοντέλου, οι τιμές αυτές είναι 1 και 0 αντίστοιχα ανεξαρτήτως SNR σε dB, όπως και πριν. Το αναμενόμενο είναι το audiovisual μοντέλο να εμπιστεύεται τον ήχο περισσότερο όσο αυξάνεται το SNR και στο test set, πράγμα που συμβαίνει, κοιτώντας τα συγκεκριμένα γραφήματα. Η τιμή της audio gate παρουσιάζει σταθερή άνοδο, ενώ αντίθετα η τιμή της video gate συνεχή πτώση. Για τιμή SNR -5dB, όπου δηλαδή ο θόρυβος είναι ισχυρότερος από το πραγματικό σήμα, η τιμή της audio gate είναι κάτω από 0.5, που σημαίνει ότι το μοντέλο έχει μάθει να εξάγει συμπεράσματα κυρίως βασισμένο στο βίντεο και όχι στον ήχο για αυτήν την τιμή θορύβου. Αντιθέτως, στην τιμή 15dB, όπου πλέον ο θόρυβος είναι πολύ ασθενέστερος από το πραγματικό σήμα, η τιμή της audio gate πλησιάζει το 0.7, κάτι που δείχνει ότι το μοντέλο αναγνωρίζει την αξιοπιστία του modality του ήχου και προσαρμόζεται ανάλογα.

\begin{figure}[H]
	\centering
	\includegraphics[width=0.7\textwidth]{test_gate_audio_vs_snr_db.png}
	\caption{Γράφημα τιμών της audio gate στο test set για διάφορες τιμές audio SNR}
	\label{test_audio_gate_vs_db}
\end{figure}

\begin{figure}[H]
	\centering
	\includegraphics[width=0.7\textwidth]{test_gate_video_vs_snr_db.png}
	\caption{Γράφημα τιμών της video gate στο test set για διάφορες τιμές audio SNR}
	\label{test_video_gate_vs_db}
\end{figure}

Τέλος, τα διαγράμματα \ref{test_loss_vs_db} και \ref{test_ter_vs_db} φανερώνουν την τελική επίδοση στο test set με τη μετρική test TER και την τιμή της αντικειμενικής συνάρτησης, για διάφορες τιμές του SNR του ηχητικού σήματος. Και πάλι, να σημειωθεί ότι οι τιμές της αντικειμενικής συνάρτησης μεταξύ του audiovisual και του αντίστοιχου audio μοντέλου δεν είναι απολύτως συγκρίσιμες, λόγω του περαιτέρω όρου auxiliary CTC. Παρατηρούμε ότι, σε κάθε περίπτωση, η αύξηση του SNR επιφέρει βελτίωση των αποτελεσμάτων, όπως ήταν και αναμενόμενο. Όσο αυξάνεται το SNR, τόσο πιο αξιόπιστο γίνεται το σήμα του ήχου, το οποίο, ακόμα και μόνο του, μπορεί να χρησιμοποιηθεί ικανοποιητικά για την πρόβλεψη φωνημάτων. Η τάση αυτή παρουσιάζεται και στις τιμές της αντικειμενικής συνάρτησης αλλά και στη μετρική test TER. Επίσης, όπως είναι αναμενόμενο, η επίδραση της αύξησης του SNR είναι μεγαλύτερη στην περίπτωση του audio μοντέλου, από ότι σε αυτήν του audiovisual μοντέλου. Αυτό οφείλεται, προφανώς, στο γεγονός ότι το μοντέλο με αποκλειστική πρόσβαση στον ήχο αντιμετωπίζει μεγαλύτερη δυσκολία στην πρόβλεψη φωνημάτων για χαμηλά SNR σε σχέση με το audiovisual μοντέλο, το οποίο διαθέτει πρόσβαση και στο modality του βίντεο (ακόμα κι αν αυτό είναι θορυβώδες, μπορεί να αξιοποιηθεί). Έτσι, η μείωση του test TER στην περίπτωση του audio μοντέλου είναι από 0.24 περίπου για -5db στο 0.16 για 15dB, ενώ για το audiovisual μοντέλο είναι από 0.18 στο 0.15 περίπου.

\begin{figure}[H]
	\centering
	\includegraphics[width=0.7\textwidth]{test_loss_vs_snr_db.png}
	\caption{Γράφημα τιμών της loss function στο test set για διάφορες τιμές audio SNR}
	\label{test_loss_vs_db}
\end{figure}

\begin{figure}[H]
	\centering
	\includegraphics[width=0.7\textwidth]{test_ter_vs_snr_db.png}
	\caption{Γράφημα τιμών του TER στο test set για διάφορες τιμές audio SNR}
	\label{test_ter_vs_db}
\end{figure}

\paragraph{Υποπείραμα 2: Αύξηση των επηρεασμένων frames του βίντεο}
\label{FramesExperiment}

Ως συμμετρικό του προηγούμενου υποπειράματος, στο 2ο υποπείραμα μελετάται η επίδραση της αύξησης των επηρεασμένων frames του βίντεο από τον θόρυβο (με τη μορφή μηδενισμού, παγώματος ή θολώματος) στην επίδοση του audiovisual μοντέλου, ενώ συγκρίνεται και με την επίδοση του μοντέλου που χρησιμοποιεί αποκλειστικά το θορυβώδες βίντεο για να προβλέψει τα φωνήματα. Το αναμενόμενο αποτέλεσμα είναι ότι όσο περισσότερα frames του βίντεο επηρεάζονται από θόρυβο, τόσο πιο αναξιόπιστο γίνεται το modality του βίντεο και άρα τόσο λιγότερο το ``εμπιστεύεται'' το μοντέλο και τόσο χειρότερες είναι οι ταξινομήσεις του σε φωνήματα. Ταυτόχρονα, το μοντέλο που χρησιμοποιεί αποκλειστικά το βίντεο για τις προβλέψεις φωνημάτων αναμένεται να πετυχαίνει χειρότερα αποτελέσματα σε σχέση με το audiovisual μοντέλο, λόγω της εκμετάλλευσης του ηχητικού σήματος. Κατά το πείραμα αυτό, η τιμή του SNR του ήχου είναι σταθερά 15 και το ηχητικό dropout παραμένει με πιθανότητα 30\% και μήκος 15 frames, ενώ η τιμή των video frames που επηρεάζονται είναι 10, 15, 20, 25 και 30 (αυτή τη φορά, η τιμή είναι σταθερή και δεν υπάρχει jitter).

Κατά την εκπαίδευση όλων των μοντέλων, πάλι χρησιμοποιήθηκε Adam optimizer με ρυθμό μάθησης 3e-4, weight decay με τιμή 1e-4 και gradient clipping με τιμή 1. Επιλέχθηκε scheduling του ρυθμού μάθησης με γραμμικό warmup για τις πρώτες 5 εποχές, μείωση του ρυθμού μάθησης μετά από 6 εποχές χωρίς βελτίωση κατά παράγοντα 0.5, με ελάχιστη τιμή το 1e-6.  Χρησιμοποιήθηκε early stopping με κριτήριο τη μετρική val TER (validation Token Error Rate) και τιμή υπομονής 12 εποχές, batch size 4, με μέγιστο αριθμό εποχών να είναι 100. Τέλος, κατά την εκπαίδευση, χρησιμοποιείται παράγοντας modality dropout 0.15 (με πιθανότητα 0.15 αφαιρείται πλήρως ένα από τα δύο modalities) για να αυξηθεί η ευρωστία του μοντέλου. Τα παρακάτω διαγράμματα περιλαμβάνουν τα αποτελέσματα του πειράματος.

Το σχήμα \ref{best_val_ter_vs_frames} δείχνει την τάση της μετρικής val TER (τη βέλτιστη τιμή της) για το audiovisual και το video μοντέλο σε συνάρτηση με τον αριθμό frames του βίντεο στα οποία έχει εφαρμοστεί θόρυβος (με τους τρόπους που αναφέρονται παραπάνω). Η αναμενόμενη εικόνα είναι η επίδοση των μοντέλων να μειώνεται όσο αυξάνονται τα επηρεασμένα frames του βίντεο, μια και το modality αυτό θα γίνεται όλο και πιο αναξιόπιστο. Επιπροσθέτως, αναμένουμε η επίδραση αυτή να είναι μεγαλύτερη για την περίπτωση του video μοντέλου, καθώς χρησιμοποιεί μονάχα το βίντεο και όχι ηχητικά στοιχεία για τις αποφάσεις του. Πράγματι, η εικόνα του σχήματος αυτού αναπαράγει τα αναμενόμενα αποτελέσματα. Η βέλτιστη τιμή της μετρικής val TER για το video μοντέλο παίρνει τιμή περίπου 0.2 για 10 επηρεασμένα frames και σταθερά ανεβαίνει, ξεπερνώντας την τιμή 0.35 για 30 επηρεασμένα frames του βίντεο. Η κατάσταση αυτή είναι αναμενόμενη, αν αναλογιστούμε ότι τα κλιπ αποτελούνται από 75 frames συνολικά, κάνοντας τα 30 από τα 75 ένα ποσοστό 40\% των εικονιζόμενων frames να έχουν επηρεαστεί. Από την άλλη, για το audiovisual μοντέλο παρατηρείται μία ελαφριά αύξηση της μετρικής όσο αυξάνει ο αριθμός των επηρεασμένων frames, από την τιμή 0.11 για 10 frames μέχρι την τιμή 0.145 για 25 frames, με μία μείωση όταν ο αριθμός των frames αυξάνεται στα 30, περίπου στην τιμή 0.13. Όπως ήταν αναμενόμενο, η επίδραση του φαινομένου υπό μελέτη είναι μικρότερη για το audiovisual μοντέλο, το οποίο σταθερά επικρατεί έναντι του visual μοντέλου. Αυτό οφείλεται, προφανώς, στην εκμετάλλευση του ηχητικού σήματος, το οποίο στις συγκεκριμένες δοκιμές είναι αρκετά ακριβές (15dB SNR, με περιπτώσεις ολικού dropout).

\begin{figure}[H]
	\centering
	\includegraphics[width=0.7\textwidth]{best_val_ter_vs_video_obstruct_frames.png}
	\caption{Γράφημα τιμών του βέλτιστου val TER για διάφορες τιμές obstructed video frames}
	\label{best_val_ter_vs_frames}
\end{figure}

Το σχήμα \ref{epoch_train_loss_across_frames} δείχνει την πορεία της αντικειμενικής συνάρτησης στο training set του audiovisual και του video μοντέλου για τις διάφορες τιμές θορυβωδών frames του βίντεο που χρησιμοποιήθηκαν. Να σημειωθεί και πάλι ότι οι τιμές της αντικειμενικής συνάρτησης του audiovisual μοντέλου και του αντίστοιχου visual μοντέλου δεν είναι απολύτως συγκρίσιμες, λόγω της έλλειψης του auxiliary CTC όρου του audio stream στην περίπτωση του μοντέλου video. Παρατηρείται σε κάθε περίπτωση ομαλή σύγκλιση, με τη βέλτιστη τελική τιμή της αντικειμενικής συνάρτησης να δείχνει να αυξάνει, όσο αυξάνονται και τα επηρεασμένα frames. Μάλιστα, για υψηλές τιμές επηρεασμένων frames, όπου το μοντέλο του βίντεο αναμένουμε να δυσκολεύεται δεδομένων των θορυβωδών συνθηκών, η τιμή της αντικειμενικής συνάρτησης του audiovisual μοντέλου φαίνεται να είναι χαμηλότερη από αυτήν του video μοντέλου, παρά της εγγενούς ανισορροπίας τους. Αυτό αποδεικνύει τη βελτίωση που παρέχει η αξιοποίηση του περαιτέρω modality. Παρά το γεγονός ότι και ο ήχος χαρακτηρίζεται από θόρυβο, το μοντέλο μαθαίνει πολύ πιο εύκολα την αναγνώριση φωνημάτων αξιοποιώντας δύο modalities έναντι ενός.

\begin{figure}[H]
	\centering
	\includegraphics[width=0.9\textwidth]{epoch_train_total_by_video_obstruct_frames.png}
	\caption{Γραφήματα τιμών train loss per epoch για διάφορες τιμές obstructed video frames}
	\label{epoch_train_loss_across_frames}
\end{figure}

Τα σχήματα \ref{epoch_audio_gate_across_frames} και \ref{epoch_video_gate_across_frames} δείχνουν τις μέσες τιμές της audio gate και της video gate, αντίστοιχα, για κάθε εποχή της εκπαίδευσης. Προφανώς, στην περίπτωση των video μοντέλων, οι τιμές αυτές είναι 0 και 1 αντίστοιχα για κάθε εποχή, καθώς δε χρησιμοποιείται εξ' ορισμού το ηχητικό σήμα. Για το audiovisual μοντέλο, αναμένουμε η μέση τιμή της audio gate να αυξάνει και, αντίστοιχα, αυτή της video gate να μειώνεται, όσο αυξάνεται ο αριθμός των frames που χαρακτηρίζονται από θόρυβο, καθώς το σήμα του βίντεο γίνεται πιο αναξιόπιστο για την πρόβλεψη φωνημάτων και άρα το μοντέλο βασίζει τις αποφάσεις του όλο και περισσότερο στο ηχητικό σήμα. Πράγματι, αυτή φαίνεται να είναι η εικόνα, μια και για τιμή επηρεασμένων frames 10 η μέση τιμή της audio gate ταλαντεύεται λίγο πάνω από το 0.5 και η αντίστοιχη της video gate λίγο κάτω από αυτήν την τιμή (το μοντέλο φαίνεται να εμπιστεύεται εξίσου τα δύο modalities), ενώ για τιμή 30 frames η μέση τιμή της audio gate ξεπερνάει το 0.6. Η διαφορά αυτή δε φαίνεται ιδιαίτερα μεγάλη, υποδεικνύοντας ότι το audiovisual μοντέλο μπορεί να εξάγει χρήσιμα συμπεράσματα από το modality του βίντεο, ακόμα και όταν αυτό είναι μεταλλαγμένο σε σημαντικό βαθμό.

\begin{figure}[H]
	\centering
	\includegraphics[width=0.9\textwidth]{epoch_val_gate_audio_by_video_obstruct_frames.png}
	\caption{Γραφήματα τιμών της audio gate στο validation set per epoch για διάφορες τιμές obstructed video frames}
	\label{epoch_audio_gate_across_frames}
\end{figure}

\begin{figure}[H]
	\centering
	\includegraphics[width=0.9\textwidth]{epoch_val_gate_video_by_video_obstruct_frames.png}
	\caption{Γραφήματα τιμών της video gate στο validation set per epoch για διάφορες τιμές obstructed video frames}
	\label{epoch_video_gate_across_frames}
\end{figure}

Τα σχήματα \ref{epoch_val_ter_across_frames} και \ref{epoch_val_loss_across_frames} δείχνουν την πορεία της μετρικής val TER και της τιμής της αντικειμενικής συνάρτησης στο validation set για τα epochs της εκπαίδευσης, με τις διάφορες τιμές αριθμού θορυβωδών frames που μελετήθηκαν. Πάλι οι τιμές της αντικειμενικής συνάρτησης του audiovisual και του αντίστοιχου video μοντέλου δεν είναι απολύτως συγκρίσιμες, για τον ίδιο λόγο που αναφέρθηκε παραπάνω. Στα συγκεκριμένα διαγράμματα, αναμενόμενη εικόνα είναι σύγκλιση, με την αύξηση των επηρεασμένων frames να επιφέρει και επιδείνωση των εκάστοτε μετρικών. Πράγματι, το σχήμα \ref{epoch_val_ter_across_frames} έχει αυτήν την εικόνα, με το video μοντέλο πάντα να πετυχαίνει χειρότερη επίδοση από το αντίστοιχο audiovisual μοντέλο, όπως είναι και λογικό. Οι επιδόσεις και των δύο στη συγκεκριμένη μετρική φαίνεται να είναι χειρότερες, όσο αυξάνεται ο αριθμός των επηρεασμένων frames, με την επίδραση αυτή να είναι μεγαλύτερη για το video μοντέλο, κάτι που ήταν αναμενόμενο. Η εικόνα του σχήματος \ref{epoch_val_ter_across_frames} φανερώνει στοιχεία overfitting για την περίπτωση του video μοντέλου, καθώς η σύγκλιση δεν είναι ομαλή και διακόπτεται από μία περίοδο αύξησης, σε κάθε περίπτωση. Να σημειωθεί ότι η συμπεριφορά αυτή μπορούσε να αποφευχθεί, χρησιμοποιώντας ως κριτήριο του early stopping το val loss και όχι το val TER. Αντιθέτως, η εικόνα του audiovisual μοντέλου δείχνει σταθερή σύγκλιση, σε τιμή η οποία μάλιστα είναι καλύτερη από την αντίστοιχη του video μοντέλου, ακόμα κι αν οι δύο αντικειμενικές συναρτήσεις παρουσιάζουν ανισορροπία. Η εκπαίδευση του audiovisual μοντέλου είναι, γενικά, πιο ευσταθής και πιο επιτυχής, αναδεικνύοντας πόσο διευκολύνεται η επίτευξη του στόχου της αναγνώρισης φωνημάτων αξιοποιώντας και στοιχεία από το ηχητικό σήμα.

\begin{figure}[H]
	\centering
	\includegraphics[width=0.9\textwidth]{epoch_val_ter_by_video_obstruct_frames.png}
	\caption{Γραφήματα τιμών val TER per epoch για διάφορες τιμές obstructed video frames}
	\label{epoch_val_ter_across_frames}
\end{figure}

\begin{figure}[H]
	\centering
	\includegraphics[width=0.9\textwidth]{epoch_val_total_by_video_obstruct_frames.png}
	\caption{Γραφήματα τιμών val loss per epoch για διάφορες τιμές obstructed video frames}
	\label{epoch_val_loss_across_frames}
\end{figure}

Το σχήμα \ref{epochs_vs_snr_db} δείχνει τον αριθμό εποχών εκπαίδευσης του audiovisual και του video μοντέλου μέχρι να πυροδοτηθεί ο μηχανισμός του early stopping, για τις διάφορες τιμές αριθμού επηρεασμένων από τον θόρυβο frames του video. Παρατηρούμε ότι, σε γενικές γραμμές, οι εποχές εκπαίδευσης του video μοντέλου τείνουν να γίνονται περισσότερες όσο αυξάνεται ο αριθμός των obstructed frames, ενώ του audiovisual μοντέλου ελαττώνονται. Η τάση αυτή δεν είναι απόλυτη και παρουσιάζει αρκετές εξαιρέσεις (για 20 frames, οι εποχές εκπαίδευσης του video model είναι περισσότερες από ότι για 25 ή 30, ενώ για 30 frames το audiovisual μοντέλο απαιτεί τις περισσότερες εποχές εκπαίδευσης σε σχέση με οποιαδήποτε άλλη περίπτωση). Το φαινόμενο αυτό οφείλεται στην ποιότητα των δεδομένων: ειδικά για το video μοντέλο, το οποίο δεν έχει πρόσβαση σε ηχητικά στοιχεία, η ποιότητα των video συσχετίζεται άμεσα με την ποιότητα και την ταχύτητα της εκπαίδευσης. Εάν τα δεδομένα είναι ιδιαίτερα θορυβώδη, το μοντέλο μπορεί να πασχίζει περισσότερο ώστε να καταλήξει σε πιθανές προβλέψεις, βασιζόμενο στα φωνήματα που έχει προβλέψει γύρω από τη θορυβώδη περιοχή του βίντεο και χτίζοντας γνώση του λεξιλογίου (ή, γενικότερα, της γλώσσας) του dataset. Έτσι, μέχρι να φτάσει στο σημείο που μπορεί να κάνει (σχετικά) ρεαλιστικές προβλέψεις για τη θορυβώδη περιοχή, απαιτείται περισσότερος χρόνος εκπαίδευσης και άρα περισσότερες εποχές. Αντίθετα, το audiovisual μοντέλο έχει τη δυνατότητα να βασιστεί όλο και περισσότερο στον ήχο, όσο το βίντεο γίνεται πιο θορυβώδες. Επομένως, αντί να περιπλέκει τη μάθησή του με περίπλοκες συσχετίσεις μεταξύ ήχου και βίντεο, μαθαίνει (ορθά) να εμπιστεύεται περισσότερο τον ήχο και ως δεύτερη επιλογή το βίντεο και μειώνει (σε γενικές γραμμές) τις απαιτούμενες εποχές για την εκπαίδευσή του.

\begin{figure}[H]
	\centering
	\includegraphics[width=0.7\textwidth]{epochs_ran_vs_video_obstruct_frames.png}
	\caption{Γράφημα αριθμού εποχών εκπαίδευσης για διάφορες τιμές obstructed video frames}
	\label{epochs_vs_frames}
\end{figure}

Τα σχήματα \ref{test_audio_gate_vs_frames} και \ref{test_video_gate_vs_frames} δείχνουν τις μέσες τιμές των audio και video gates αντίστοιχα στο test set για το audiovisual και το video μοντέλο, για τις τιμές obstructed frames που δοκιμάστηκαν. Προφανώς οι τιμές αυτές είναι 0 και 1 αντίστοιχα στην περίπτωση του video μοντέλου. Για το audiovisual μοντέλο, είναι αναμενόμενο να αυξάνεται η τιμή της audio gate και αντίστοιχα να μειώνεται αυτή της video gate όσο αυξάνεται ο αριθμός των obstructed frames, καθώς το μοντέλο μαθαίνει να εμπιστεύεται το πιο αξιόπιστο σήμα, το οποίο στη συγκεκριμένη περίπτωση είναι ο ήχος. Πράγματι, αυτή είναι η εικόνα που παρουσιάζουν τα συγκεκριμένα διαγράμματα, με διαρκή αύξηση της μέσης τιμής της audio gate και αντίστοιχη μείωση της μέσης τιμής της video gate, όσο αυξάνεται ο αριθμός των επηρεασμένων από τον θόρυβο frames. Το μοντέλο μαθαίνει να δίνει όλο και περισσότερη έμφαση στα συμπεράσματα που εξάγει από τον ήχο και να χρησιμοποιεί συμπληρωματικά το βίντεο, όσο το βίντεο γίνεται πιο θορυβώδες και άρα πιο αναξιόπιστο.

\begin{figure}[H]
	\centering
	\includegraphics[width=0.7\textwidth]{test_gate_audio_vs_video_obstruct_frames.png}
	\caption{Γράφημα τιμών της audio gate στο test set για διάφορες τιμές obstructed video frames}
	\label{test_audio_gate_vs_frames}
\end{figure}

\begin{figure}[H]
	\centering
	\includegraphics[width=0.7\textwidth]{test_gate_video_vs_video_obstruct_frames.png}
	\caption{Γράφημα τιμών της video gate στο test set για διάφορες obstructed video frames}
	\label{test_video_gate_vs_frames}
\end{figure}

Τέλος, τα σχήματα \ref{test_loss_vs_frames} και \ref{test_ter_vs_frames} απεικονίζουν τις τιμές της αντικειμενικής συνάρτησης στο test set και της μετρικής test TER, για το audiovisual και video μοντέλο, στις συνθήκες που δοκιμάστηκαν. Από τα παραπάνω, αναμενόμενο αποτέλεσμα είναι η επίδοση των μοντέλων να επιδεινώνεται όσο αυξάνεται ο αριθμός των θορυβωδών frames, με το φαινόμενο αυτό να γίνεται πιο αντιληπτό στο video μοντέλο. Αναμένουμε, επίσης, το audiovisual μοντέλο σε κάθε περίπτωση να επιφέρει βελτιωμένα αποτελέσματα σε σχέση με το video μοντέλο, εκμεταλλευόμενο την πρόσβασή του στο ηχητικό σήμα. Η εικόνα αυτή φαίνεται πράγματι να παρουσιάζεται στα αποτελέσματα. Παρά το γεγονός ότι οι αντικειμενικές συναρτήσεις των μοντέλων δεν είναι απολύτως συγκρίσιμες (για τους λόγους που αναφέρθηκαν παραπάνω), η τιμή της αντικειμενικής συνάρτησης του audiovisual μοντέλου είναι σε κάθε περίπτωση μικρότερη από την αντίστοιχη τιμή του video μοντέλου, αποδεικνύοντας τη διαφορά στην ποιότητα της μάθησης που επιτυγχάνουν τα δύο μοντέλα στο θορυβώδες περιβάλλον. Η τιμή της αντικειμενικής συνάρτησης και η μετρική test TER παρουσιάζουν μία γενικά αυξητική τάση και για τα δύο μοντέλα, όσο αυξάνεται ο αριθμός των επηρεασμένων frames, τάση η οποία, όπως ήταν αναμενόμενο, είναι ενισχυμένη για το video μοντέλο. Τα αποτελέσματα αυτά αναδεικνύουν τη διευκόλυνση του προβλήματος που παρέχεται μέσω της πρόσβασης του, έστω και θορυβώδους, ηχητικού σήματος.

\begin{figure}[H]
	\centering
	\includegraphics[width=0.7\textwidth]{test_loss_vs_video_obstruct_frames.png}
	\caption{Γράφημα τιμών της loss function στο test set για διάφορες τιμές obstructed video frames}
	\label{test_loss_vs_frames}
\end{figure}

\begin{figure}[H]
	\centering
	\includegraphics[width=0.7\textwidth]{test_ter_vs_video_obstruct_frames.png}
	\caption{Γράφημα τιμών του TER στο test set για διάφορες τιμές audio SNR}
	\label{test_ter_vs_frames}
\end{figure}